
\part{The problem and the state of knowledge}

%%==================================================================
\section{Definitions, history and equivalent forms}\label{sec:history}
%%==================================================================

\subsection{Time, nondeterminism and reductions}
Hartmanis and Stearns \cite{HS65} measured the difficulty of a problem by the
running time of the best Turing machine that decides it and proved the first
hierarchy theorem: for time-constructible $t$, more time decides more
languages. A nondeterministic machine may branch at each step and accepts when
some branch accepts. The class $\NP=\bigcup_k\NTIME(n^k)$ is the set of
languages $\mathcal L$ for which there are a polynomial $p$ and a polynomial-time
relation $R$ with
\[
x\in\mathcal L\iff \exists y\in\{0,1\}^{p(|x|)}\ R(x,y).
\]
The string $y$ is a \emph{certificate}. A polynomial-time \emph{reduction} from
$\mathcal L$ to $\mathcal L'$ is a polynomial-time map $g$ with $x\in\mathcal
L\iff g(x)\in\mathcal L'$, and $\mathcal L'$ is $\NP$-complete if it lies in
$\NP$ and every language in $\NP$ reduces to it.

Cook \cite{Coo71} and, independently, Levin \cite{Lev73} proved that the
satisfiability of propositional formulas is $\NP$-complete.\footnote{Levin's
note treats search problems (``universal sequential search problems'') and
gives six complete ones, among them tiling and set cover; Cook works with
decision problems and Turing reductions. Karp's reductions are the many-one
reductions used today.} The proof writes the computation of a nondeterministic
machine on $x$ as a table of $p(|x|)^2$ cells and expresses its correctness by
clauses that each read a bounded window of the table. Karp \cite{Kar72} reduced
satisfiability to twenty-one problems of combinatorics and optimisation,
including clique, vertex cover, Hamiltonian cycle, three-dimensional matching,
knapsack and zero-one integer programming, and Garey and Johnson \cite{GJ79}
catalogued several hundred more. Consequently $\Pc=\NP$ if and only if any one
of these problems has a polynomial-time algorithm.

Edmonds \cite{Edm65} had argued, in his paper on matchings, that polynomial time
is the right notion of a ``good'' algorithm. G\"odel asked von Neumann in 1956 whether
the existence of proofs of length $n$ for a given formula can be decided in time
linear or quadratic in $n$, which is the bounded form of the same question; the
letter is translated and discussed by Sipser \cite{Sip92}.

\subsection{What is known to separate}
Diagonalisation separates classes defined by the same resource at different
amounts. The deterministic hierarchy of Hartmanis and Stearns, and the
nondeterministic hierarchies of Cook \cite{Coo73}, Seiferas, Fischer and Meyer
\cite{SFM78} and \v{Z}\'ak \cite{Zak83}, show for instance
$\NTIME(2^n)\ne\NTIME(2^{n/2})$; \v{Z}\'ak's form, which separates unary
languages, is the one used in Section~\ref{sec:lower}. Between different
resources very little is separated. Paul, Pippenger, Szemer\'edi and Trotter
\cite{PPST83} proved $\DTIME(n)\ne\NTIME(n)$ for multitape machines, a statement
that does not extend to $n^k$. Hopcroft, Paul and Valiant \cite{HPV77} showed
that time $t$ can be simulated in space $t/\log t$, and Williams \cite{Wil25}
improved this to space $\sqrt{t\log t}$. Fortnow, Lipton, van Melkebeek and
Viglas \cite{FLMV05} proved that satisfiability cannot be decided
simultaneously in time $n^c$ and space $n^{o(1)}$ for any $c$ below the golden
ratio.

For space, nondeterminism costs at most a square: Savitch \cite{Sav70} proved
that nondeterministic space $s$ is contained in deterministic space $s^2$, so
$\mathsf{PSPACE}=\mathsf{NPSPACE}$. No analogue is known for time; the best
deterministic simulation of $\NTIME(t)$ takes time $2^{O(t)}$.

\subsection{Structure inside $\NP$}
If $\Pc\ne\NP$, Ladner \cite{Lad75} showed that $\NP$ contains problems that are
neither in $\Pc$ nor $\NP$-complete. Some natural problems were long thought to
be of this kind; primality was placed in $\Pc$ by Agrawal, Kayal and Saxena
\cite{AKS04}. Counting is much stronger than deciding: Valiant \cite{Val79b}
proved the permanent of a zero-one matrix $\#\Pc$-complete although the
corresponding decision problem, existence of a perfect matching, is in $\Pc$,
and Toda \cite{Tod91} showed that the whole polynomial hierarchy reduces to
$\#\Pc$. The PCP theorem of Arora and Safra \cite{AS98} and of Arora, Lund,
Motwani, Sudan and Szegedy \cite{ALMSS98}, with Dinur's combinatorial proof
\cite{Din07}, shows that every language in $\NP$ has certificates that a
randomised verifier checks by reading a constant number of bits; equivalently,
approximating the maximum fraction of satisfiable clauses of a $3$-CNF within
some constant factor is already $\NP$-hard. Impagliazzo \cite{Imp95} described
five possible worlds, from $\Pc=\NP$ to the existence of public-key
cryptography, that are consistent with present knowledge.

\subsection{Circuits}
A Boolean circuit with $n$ inputs is a directed acyclic graph whose sources are
the inputs and whose other vertices, the gates, compute Boolean functions of
their predecessors. The \emph{size} of a circuit over the binary basis is its
number of gates; for gates of larger fan-in the right measure is the number of
\emph{wires}, the sum of the fan-ins. A language $\mathcal L$ has circuits of
size $s(n)$ if for each $n$ some circuit of that size decides $\mathcal L$ on
inputs of length $n$; the class $\Ppoly$ consists of the languages with
circuits of polynomial size. Since $\Pc\subseteq\Ppoly$, a proof that some
language in $\NP$ lies outside $\Ppoly$ would prove $\Pc\ne\NP$. The circuit
form is attractive because it removes uniformity, so that only combinatorics
remains, and because $\NP\not\subseteq\Ppoly$ is believed to be true for
satisfiability itself. Kannan \cite{Kan82} proved that $\Sigma_2\Pc$ is not
contained in $\SIZE(n^k)$ for any fixed $k$, and Santhanam \cite{San09}
proved the same for Merlin--Arthur languages with one bit of advice.

\subsection{The Millennium Problem}
The Clay Mathematics Institute listed the question among its seven Millennium
Prize Problems in 2000 \cite{Clay}; the official problem description is by
Cook. Fortnow \cite{For09} and Aaronson \cite{Aar16} survey the evidence,
Wigderson \cite{Wig19} places it within mathematics as a whole, and the
textbooks of Arora and Barak \cite{AB09} and of Jukna \cite{Juk12} give the
technical background used below. Gasarch's polls \cite{Gas19} record a
steadily growing majority of researchers who expect $\Pc\ne\NP$.

%%==================================================================
\section{Three barriers}\label{sec:barriers}
%%==================================================================

Three theorems show that the methods that settle other questions cannot settle
this one by themselves. Each barrier is a statement about proofs, and each has
been circumvented in some restricted setting.

\subsection{Relativization}
Baker, Gill and Solovay \cite{BGS75} constructed oracles $A$ and $B$ with
$\Pc^A=\NP^A$ and $\Pc^B\ne\NP^B$. Any argument that treats machines as black
boxes, such as diagonalisation and simulation, gives the same answer relative
to every oracle and therefore cannot decide the question. The hierarchy
theorems relativize; the Cook--Levin theorem, which uses the local structure of
computations, does not relativize in its strong forms, and the arguments of
Section~\ref{sec:algorithms} use it through local reductions.

\subsection{Natural proofs}
Razborov and Rudich \cite{RR97} observed that the known circuit lower bounds
proceed by a property $\Gamma$ of truth tables that is \emph{useful} (no
function with the property has small circuits), \emph{large} (a fraction
$2^{-O(n)}$ of all functions has it) and \emph{constructive} (it can be
decided in time polynomial in $2^n$ from the truth table). They proved that a
natural property useful against $\Ppoly$ would break pseudorandom functions
that are secure against circuits of size $2^{n^\epsilon}$, whose existence
follows from standard cryptographic assumptions. Arguments by diagonalisation
are not natural, since they single out one function; Williams \cite{Wil16}
showed that $\NEXP\not\subseteq\mathcal C$ is in fact equivalent to the
existence of a property useful against $\mathcal C$ that is constructive in a
weak, nondeterministic sense, which locates the barrier exactly.

\subsection{Algebrization}
Aaronson and Wigderson \cite{AW09} extended relativization to algebraic
oracles, which also capture the arithmetisation used in interactive proofs.
They proved that every known non-relativizing technique of that kind
algebrizes, and that resolving $\Pc$ versus $\NP$ needs non-algebrizing
methods.

\subsection{Locality}
A fourth, newer observation concerns hardness magnification. Oliveira and
Santhanam \cite{OS18}, and earlier Allender and Kouck\'y \cite{AK10}, found
problems for which a slightly superlinear lower bound would imply a
superpolynomial one. Chen, Hirahara, Oliveira, Pich, Rajgopal and Santhanam
\cite{CHOPRS22} proved that the known slightly superlinear lower bounds hold
also for circuits with small oracle gates of unbounded function, while the
magnified statements fail for such circuits; so the methods that prove the weak
bounds cannot prove the magnified ones. This is the \emph{locality barrier}.

\begin{remark}\label{rem:barrierhere}
The lower bound of Theorem~B is proved by diagonalisation against a
nondeterministic class combined with a satisfiability algorithm. It names a
single language, so it is not natural, and it uses the local structure of
computations through the reduction of Lemma~\ref{lem:jmv}.
Section~\ref{sec:limits} discusses its relation to the locality barrier.
\end{remark}

%%==================================================================
\section{Circuit lower bounds known today}\label{sec:circuits}
%%==================================================================

Figure~\ref{fig:landscape} arranges the results of this section by model and
strength.

\begin{figure}[t]
\centering
\includegraphics[width=0.92\textwidth]{fig_landscape.png}
\caption{Explicit circuit lower bounds by model. The height of a bar is the
order of growth of the best bound, on a scale from linear ($cn$) to exponential
($2^{n^\epsilon}$); ``barely superlinear'' means $n\cdot\omega(1)$ but
$n^{1+o(1)}$. For circuits of unbounded depth over a complete basis (orange),
only linear bounds were known for explicit languages before \cite{OAI26};
\cite[Corollary 1.4]{OAI26} and Theorem~B give $cn\sqrt{\log n}$. The bounds for
arbitrary gates of bounded depth are for operators with $n$ outputs.}\label{fig:landscape}
\end{figure}

\subsection{General circuits}
Shannon's counting \cite{Sha49} shows that almost every function needs about
$2^n/n$ gates, and Lupanov's construction (see \cite[Chapter 1]{Juk12})
matches this. For explicit functions the only method for unbounded depth has
been gate elimination: fixing an input variable removes a few gates, and
repeating the step $n$ times gives a linear bound. Schnorr \cite{Sch80} and
Blum \cite{Blu84} pushed the constant to three (Schnorr announced $3n$, and
Blum, noting a gap in one lemma, proved $3n-o(n)$ for a modification of Paul's
function), Find, Golovnev, Hirsch and
Kulikov \cite{FGHK16} to $3+\frac1{86}$, and Li and Yang \cite{LY22} to
$3.1n-o(n)$ for a function in $\Pc$.\footnote{Over the de Morgan basis, where
negations are free and the gates are AND and OR, Iwama and Morizumi
\cite{IM02} proved $5n-o(n)$.} Golovnev, Hirsch, Knop and Kulikov
\cite{GHKK18} showed that natural forms of gate elimination cannot prove much
more than linear bounds, and Golovnev, Kulikov, Smal and Tamaki \cite{GKST18}
turned gate elimination into satisfiability algorithms for linear-size
circuits.

\subsection{Formulas}
A formula is a circuit whose gates have fan-out one. Random restrictions
shrink de Morgan formulas: H\aa stad \cite{Has98} proved shrinkage exponent
two, so Andreev's function requires formulas of size $n^{3-o(1)}$, and Tal
\cite{Tal14} sharpened the lower-order terms to a polylogarithmic factor.
Karchmer and Wigderson \cite{KW90} identified formula depth with the
communication complexity of a game; with it they proved that monotone formulas
for connectivity need depth $\Omega(\log^2 n)$, and Raz and McKenzie
\cite{RM99} separated the levels of the monotone $\mathsf{NC}$ hierarchy.

\subsection{Bounded depth}
Ajtai \cite{Ajt83} and Furst, Saxe and Sipser \cite{FSS84} proved that parity is
not in $\AC$, the class of polynomial-size circuits of constant depth with
unbounded fan-in AND and OR gates. Yao \cite{Yao85} and H\aa stad \cite{Has86}
made the bound exponential: depth $d$ requires size $2^{\Omega(n^{1/(d-1)})}$.
Rossman \cite{Ros08} proved that detecting a $k$-clique needs size
$\omega(n^{k/4})$ at every constant depth. With parity gates added, Razborov
\cite{Raz87} and Smolensky \cite{Smo87} proved exponential lower bounds for
majority, and for counting modulo $q$ in circuits with gates modulo a prime
$p\nmid q$.

\subsection{Monotone circuits}
Razborov's method of approximations gave superpolynomial lower bounds for
clique in monotone circuits, made exponential by Alon and Boppana \cite{AB87}
(see \cite[Chapter 9]{Juk12}). Tardos \cite{Tar88} showed that the gap between
monotone and general circuits can be exponential, so monotone bounds do not
transfer.

\subsection{$\ACC$ and the algorithmic method}
For circuits of constant depth with gates counting modulo arbitrary $m$, the
class $\ACC$, no lower bound followed from the methods above. Williams
\cite{Wil14} proved $\NEXP\not\subseteq\ACC$ by a satisfiability algorithm for
$\ACC$ circuits that beats exhaustive search, combined with the argument of
\cite{Wil13} described in Section~\ref{sec:algorithms}. Murray and Williams
\cite{MW20} extended it to nondeterministic quasi-polynomial time, and
Williams \cite{Wil14b} to $\ACC$ circuits with a layer of threshold gates at the
bottom.

\subsection{High classes}
Exponential classes with alternation reach maximum circuit size. Chen,
Hirahara and Ren \cite{CHR24} proved that a language in symmetric exponential
time $\mathsf{S_2E}$ with one bit of advice requires circuits of size about
$2^n/n$, and Li \cite{Li24} removed the advice; this class lies between
$\ENP$ and $\mathsf{ZPE}^{\NP}$ by Cai's theorem \cite{Cai07}. For $\ENP$
itself, whether every language has circuits of linear size was open until
\cite[Corollary C.1]{OAI26}.

\subsection{Threshold circuits}
A threshold gate outputs $[\sum_iw_iy_i\ge\theta]$. Polynomial-size threshold
circuits of constant depth form $\TC$, a class that contains all of the
arithmetic of integers and for which no superpolynomial lower bound is known
even at depth two for functions in $\NEXP$. The known results concern size near
linear. Hajnal, Maass, Pudl\'ak and Szegedy \cite{HMPST93} separated depth two
from depth three for majority circuits. Impagliazzo, Paturi and Saks
\cite{IPS97} proved that parity needs $n^{1+c\theta^{-d}}$ wires in threshold
circuits of depth $d$, for constants $c>0$ and $\theta>1$. Kane and Williams
\cite{KW16} proved superlinear gate and superquadratic wire bounds for depth
two. Chen, Santhanam and Srinivasan \cite{CSS18} proved average-case lower
bounds and satisfiability algorithms for threshold circuits of depth $d$ with
$n^{1+\epsilon_d}$ wires, and Impagliazzo, Paturi and Schneider \cite{IPS13}
gave a satisfiability algorithm for sparse depth-two threshold circuits.
Every one of these bounds degrades with the depth and gives nothing at
unbounded depth.

\subsection{Arbitrary gates}
When gates may compute any function, a single gate of fan-in $n$ computes every
Boolean function, so lower bounds concern either many outputs or bounded fan-in.
For operators with $n$ outputs, Valiant's graph-theoretic method \cite{Val77},
Pudl\'ak's analysis of communication in bounded depth \cite{Pud94}, the bounds
of Radhakrishnan and Ta-Shma for depth-two superconcentrators \cite{RTS00}, and
the explicit operators of Cherukhin \cite{Che11} and Jukna \cite{Juk10} give
superlinear wire bounds that hold only at bounded depth. Theorem~B is a bound of
the second kind: its table gates have fan-in at most $N^{3/4}$ and its depth is
unrestricted.

\subsection{Algebraic circuits}
Valiant \cite{Val79a} introduced the algebraic classes $\mathsf{VP}$ and
$\mathsf{VNP}$ and showed that the permanent is complete for $\mathsf{VNP}$; the
conjecture $\mathsf{VP}\ne\mathsf{VNP}$ is the algebraic analogue of $\Pc\ne\NP$.
Limaye, Srinivasan and Tavenas \cite{LST21} proved the first superpolynomial
lower bounds against algebraic circuits of every constant depth. Kabanets and
Impagliazzo \cite{KI04} showed that derandomising polynomial identity testing
implies lower bounds of this kind, which ties the algebraic question to
pseudorandomness.

\subsection{Range avoidance}
Given a circuit $C:\{0,1\}^n\to\{0,1\}^{n+1}$, find a string outside its range.
Korten \cite{Kor21} showed that this search problem captures explicit
constructions, among them hard truth tables; Ren, Santhanam and Wang
\cite{RSW22} studied its complexity, and Ilango, Li and Williams \cite{ILW23}
connected it to indistinguishability obfuscation. The lower bounds of
\cite{CHR24,Li24} come from an algorithm for range avoidance with an $\NP$
oracle.

%%==================================================================
\section{Algorithms, the exponential time hypothesis and its refutation}\label{sec:algorithms}
%%==================================================================

\subsection{Exponential algorithms for satisfiability}
Exhaustive search decides satisfiability of a formula with $n$ variables in
time $2^n$ times a polynomial. For $k$-CNF formulas much better is known.
Sch\"oning's random walk \cite{Sch99} runs in time $(2-2/k)^n$; the algorithm of
Paturi, Pudl\'ak, Saks and Zane \cite{PPSZ05} improves this, and Hertli
\cite{Her14} proved that its bounds for formulas with a unique solution hold in
general, about $1.308^n$ for $3$-SAT. Williams \cite{Wil05} solved maximum
two-constraint satisfaction in time $2^{\omega n/3}$ by fast matrix
multiplication, where $\omega<2.38$ is the exponent of matrix multiplication.

\subsection{The hypotheses}
Impagliazzo and Paturi \cite{IP01} formulated the \emph{exponential time
hypothesis} (ETH): there is $\delta>0$ such that $3$-SAT cannot be decided in
time $2^{\delta n}$. The \emph{strong} hypothesis (SETH) asserts that for every
$\epsilon>0$ some $k$ makes $k$-SAT require time $2^{(1-\epsilon)n}$.
Impagliazzo, Paturi and Zane \cite{IPZ01} proved the sparsification lemma, by
which ETH is equivalent to the same statement for formulas with $O(n)$
clauses, and Calabro, Impagliazzo and Paturi \cite{CIP06,CIP09} related clause
width to clause density and extended the hypotheses to circuits of small depth.
A large part of fine-grained complexity derives conditional lower bounds from
these hypotheses.

\subsection{Williams' algorithmic method}
Williams \cite{Wil13} proved that a satisfiability algorithm for a circuit
class that is only slightly faster than exhaustive search yields a lower bound
for that class. The argument runs as follows (Figure~\ref{fig:pipeline}).
Suppose that every language in $\NEXP$ has small circuits in the class. By the
easy-witness lemma of Impagliazzo, Kabanets and Wigderson \cite{IKW02}, every
satisfiable instance then has a witness whose bits are computed by a small
circuit. A nondeterministic machine guesses such a circuit, composes it with a
succinct description of the instance, and verifies the witness by running the
fast satisfiability algorithm on the composed circuit. This solves every
language of $\NTIME(2^n)$ in nondeterministic time below $2^n$, contradicting
the nondeterministic hierarchy theorem. The survey \cite{Wil14c} describes the
many variants. The succinct description is a \emph{local reduction}: Jahanjou,
Miles and Viola \cite{JMV15} showed that the clauses of the Cook--Levin formula
for a computation of length $T$ can be indexed by $\log T+O(\log\log T)$ bits
in such a way that each clause is computed from its index by a circuit of
quasi-linear size whose output bits each depend on few index bits.

\begin{figure}[t]
\centering
\includegraphics[width=0.92\textwidth]{fig_pipeline.png}
\caption{Williams' argument as used in Section~\ref{sec:lower}. A unary
language that is hard for nondeterministic time $2^{n/2}$ is reduced locally to
a succinct $3$-CNF; a small circuit for the witness language $\Wl$ is guessed,
three copies of it read the variables of an indexed clause, and the test
circuit is checked by the counting algorithm of Theorem~A.}\label{fig:pipeline}
\end{figure}

Santhanam \cite{San10} gave satisfiability algorithms for formulas, Seto and
Tamaki \cite{ST13} for formulas over the full binary basis, and Chen, Kabanets,
Kolokolova, Shaltiel and Zuckerman \cite{CKKSZ15} showed how lower-bound proofs
by restriction yield such algorithms in general.

\subsection{The refutation of ETH}
In October 2026 \oai{} released the preprint \cite{OAI26}, whose main theorem
counts the satisfying assignments of a $k$-CNF with $n$ variables and $m$
clauses in time
\[
(n+m)^{O_k(1)}\,2^{O_k(n\log\log n/\sqrt{\log n})},
\]
and in time $2^{O(n/\sqrt{\log n})}$ when $m=O(n)$. This contradicts ETH, and
with the sparsification lemma it removes the basis of many conditional lower
bounds. The fixed-density bound rests on a theorem about tensor networks,
\cite[Theorem D.1]{OAI26}, restated below as Theorem~\ref{thm:contraction}: a
bipartite Boolean tensor network with $O(s)$ edges and degrees at most
$s^{9/10}$ can be contracted in time $2^{O(s/\sqrt{\log s})}$. Its proof
deforms the identification of the wires by a Cauchy matrix, expands the
interactions between groups of vertices in variables whose two-adic valuations
grow with their weight, and recovers the value at the undeformed point by
Hermite interpolation with tapered multiplicities; Appendix~\ref{app:contraction}
gives the full argument. Contraction of tensor networks had been used for the
simulation of quantum circuits by Markov and Shi \cite{MS08}, with running time
exponential in the treewidth; the theorem of \oai{} needs no bound on the
treewidth.

Combining the fixed-density algorithm with Williams' method, \cite[Corollary
1.4]{OAI26} proves that some language in $\ENP$ needs more than $cn\sqrt{\log
n}$ gates of fan-in two for infinitely many $n$. This was the first superlinear
lower bound of any kind for $\ENP$. Weighted counting, Max-Cut profiles and
exact optimisation follow from the same algorithm \cite{BhaHex26,Bha26eth}.

\subsection{What the refutation does not give}
A subexponential algorithm for $3$-SAT does not place $\NP$ in $\Pc$: the time
$2^{O(n/\sqrt{\log n})}$ is still superpolynomial. Nor does it give a
superpolynomial circuit lower bound, since Williams' method needs an algorithm
whose saving over $2^n$ is measured against the number $n$ of inputs of
circuits of polynomial size, and the algorithms of \cite{OAI26} are fast only
when the circuit has $O(n\sqrt{\log n})$ wires.

%%==================================================================
\section{Other programmes}\label{sec:otherroutes}
%%==================================================================

\subsection{Geometric complexity theory}
Mulmuley and Sohoni \cite{MS01} proposed to separate the permanent from the
determinant, and so $\mathsf{VNP}$ from $\mathsf{VP}$, by representation
theory: the orbit closure of the padded permanent would not be contained in
that of the determinant, and an \emph{obstruction}, an irreducible
representation occurring in the coordinate ring of one closure and not of the
other, would witness this. Mulmuley \cite{Mul12} describes the programme.
B\"urgisser, Ikenmeyer and Panova \cite{BIP19} proved that occurrence
obstructions, those detected by the mere presence of a representation, do not
exist in the required range; obstructions by multiplicities remain possible.

\subsection{Proof complexity}
Cook and Reckhow \cite{CR79} showed that $\NP\ne\coNP$ if and only if no
propositional proof system has polynomial-size proofs of all tautologies; since
$\Pc=\NP$ implies $\NP=\coNP$, superpolynomial lower bounds for every proof
system would separate $\Pc$ from $\NP$. Haken \cite{Hak85} proved that
resolution needs exponential size for the pigeonhole principle, and
Ben-Sasson and Wigderson \cite{BSW01} related resolution size to width.
Kraj\'i\v{c}ek's monograph \cite{Kra19} surveys the field; for Frege systems
no superpolynomial bound is known.

\subsection{Meta-complexity}
The minimum circuit size problem MCSP asks, given a truth table and a number
$s$, whether the function has a circuit of size $s$. Kabanets and Cai
\cite{KC00} began its modern study; it is in $\NP$, and its $\NP$-hardness is open.
Hirahara \cite{Hir18} gave a non-black-box reduction from worst case to average
case for approximating it, and \cite{Hir22} proved $\NP$-hardness of a partial
version. These questions concern the complexity of complexity itself and are
closely tied to natural proofs and to magnification.

\subsection{Hardness versus randomness}
Nisan and Wigderson \cite{NW94} built pseudorandom generators from hard
functions, and Impagliazzo and Wigderson \cite{IW97} proved $\mathsf{P}=
\mathsf{BPP}$ if some language in $\Ec$ requires circuits of size $2^{\Omega(n)}$.
Conversely, derandomization implies lower bounds \cite{KI04,Wil16}. A lower
bound of the strength of Theorem~B is far too weak for derandomization.

\subsection{Magnification}
The results of \cite{AK10,OS18,CHOPRS22} identify problems, many of them
variants of MCSP, for which a lower bound of the form $n^{1+\epsilon}$ in some
model implies a superpolynomial bound in a stronger model. The locality
barrier of Section~\ref{sec:barriers} explains why the known methods have not
reached the magnified thresholds.

%%==================================================================
\section{Repositories and the literature catalogue}\label{sec:repositories}
%%==================================================================

\subsection{Collections of claims}
Woeginger's page \cite{Woe16}, maintained until 2016, lists the announced
solutions that came to his attention, with entries dated from 1986 to 2016.
Parsing the page gives 114 entries with a stated verdict: 62 claim $\Pc=\NP$, 49 claim $\Pc\ne\NP$, two
claim that the question is unprovable and one that it is undecidable
(Figure~\ref{fig:claims}). None has been accepted. The first entry already
illustrates a lesson that became a theorem: polynomial linear programs for the
travelling salesman problem were refuted by Yannakakis \cite{Yan91}, who
proved that every symmetric linear program expressing the problem has
exponential size.

\begin{figure}[t]
\centering
\includegraphics[width=0.9\textwidth]{fig_claims.png}
\caption{Announced solutions of the $\Pc$ versus $\NP$ problem listed by
Woeginger \cite{Woe16}, by year and verdict (114 dated entries). None of them
is accepted.}\label{fig:claims}
\end{figure}

\subsection{Archives}
The Electronic Colloquium on Computational Complexity \cite{ECCC} and the
Complexity Zoo \cite{Zoo}, a catalogue of several hundred complexity
classes and their known inclusions, are the standard reference points. The
preprint \cite{OAI26} and its Lean formalisation are in the public repository
of \oai{}. Archived code accompanies the author's related papers on counting
\cite{Bha26eth}, on Weil classes \cite{Bha26weil} and on Tate--Shafarevich
groups \cite{Bha26bsd}; the programs that check the finite identities of this
paper are archived in the same way (Declarations).

\subsection{A catalogue of arXiv papers}
To measure the field, nineteen keyword queries were run against the arXiv
interface, one for each topic of Sections~\ref{sec:history}
to~\ref{sec:otherroutes}: $\Pc$ versus $\NP$, circuit lower bounds, natural
proofs, relativization, algebrization, geometric complexity theory, hardness
magnification, MCSP, meta-complexity, proof complexity, the exponential time
hypotheses, threshold circuits, satisfiability algorithms for circuits, formula
lower bounds, $\ACC$, $\mathsf{VP}$ versus $\mathsf{VNP}$, range avoidance, PCPs
and tensor network contraction. The union contains 6,480 papers, from 1993 to
the end of October 2026.\footnote{The queries combine words in the abstract
with category filters, and some are broad: the words ``meta'' and
``complexity'' also match papers in machine learning. Of the 6,480 records,
1,452 have primary category \texttt{cs.CC} and 2,617 lie in the core categories
of theoretical computer science and discrete mathematics; the figure and the
counts below use the core set. The catalogue, the queries and the program
that produced them are in the archive.} Figure~\ref{fig:catalogue} shows the
number of papers in the core categories by year and topic. The exponential
time hypotheses (657 papers), the $\Pc$ versus $\NP$ question itself (436),
circuit lower bounds (407) and proof complexity (332) are the largest topics;
the yearly number of papers grew from about thirty in 2003 to more than three
hundred in each of 2025 and 2026 in the core set.

\begin{figure}[t]
\centering
\includegraphics[width=0.95\textwidth]{fig_catalogue.png}
\caption{Papers of the arXiv catalogue in the core categories, by year of first
version and by topic (2,617 papers; a paper matching several topics is counted
once for each).}\label{fig:catalogue}
\end{figure}

\subsection{Where the problem stands}
No proof that $\Pc\ne\NP$ is known, and no known method is expected to give
one without new ideas, since each of the three barriers excludes a family of
techniques. The strongest explicit lower bounds for unrestricted circuits are
linear for functions in $\NP$ and, since \cite{OAI26}, of order $n\sqrt{\log
n}$ for a function in $\ENP$. Part~II extends the latter bound to circuits of
threshold, modular, symmetric and large table gates.

\part{New results}
